\centerline{\bf \Huge Фукнция Эйлера}

\textbf{Определение.} Пусть $n$ - натуральное число. Количество чисел, не превосходящих $n$ и взаимно простых с $n$, обозначается через $\varphi(n)$ («фи от эн»). Функция $\varphi$ называется \textit{функцией Эйлера}. Например, $\varphi(6)=2$, потому что среди чисел от 1 до 6 только числа 1 и 5 взаимно просты с 6 .

\problem{1}
Пусть $p$ - простое число. Чему равно $\varphi(p)$? А $\varphi\left(p^a\right)$?

\problem{2}
\textbf{Мультипликативность функции Эйлера.}
В таблицу $a \times b$ выписаны все последовательные числа подряд начиная с 1 , причем $(a, b)=1$
а) Сколько в таблице чисел, взаимно простых с $b$ ?
б) Сколько в каждом столбце чисел, взаимно простых с $a$ ?
в) Используя таблицу докажите, что $\varphi(a b)=\varphi(a) \varphi(b)$.

\problem{3}
Как вычислить $\varphi(n)$ для $n=p_1^{\alpha_1} \cdot \ldots \cdot p_k^{\alpha_k}$ ?


\problem{4}
Отдельные пункты не принимаются. Рассмотрим все дроби вида $a / n$, где $a \leqslant n$. Сократим их все.

(a) Какие числа являются знаменателями получившихся дробей?

(б) У скольких дробей знаменателем является число $d$ ?

(в) Сформулируйте и докажите соответствующее равенство.

\problem{5}
Отдельные пункты не принимаются.

(a) Пусть числа $a$ и $b$ взаимно просты с $n$. Докажите, что $a b$ также взаимно просто c $n$.

(б) Пусть числа $a, b$ и $c$ взаимно просты $c n$, а также числа $a$ и $b$ дают разные остатки при делении на $n$. Докажите, что числа $a c$ и $b c$ также дают разные остатки при делении на $n$.

(в) Пусть числа $a$ и $b$ взаимно просты с п. Докажите, что найдётся такое число с, что $a c \equiv b(\bmod\ n)$.

(г) Пусть $\mathbb{Z}_n^*$ - множество из всех чисел, не превосходящих $n$ и взаимно простых с $n$, и $a$ взаимно просто с $n$. Умножим все числа множества $S$ на $a$. Как изменится остаток от деления произведения всех чисел множества $\mathbb{Z}_n^*$ на $n$ ?

\newpage

\hspace{3mm}

(д) \textbf{Теорема Эйлера.} Докажите, что если $a$ взаимно просто с $n$, то $a^{\varphi(n)} \equiv 1(\bmod\ n)$.

\problem{6}
Пусть $a$ взаимно просто с 1001. Докажите, что $a^{3000}-1$ делится на 1001.

\problem{7}
Существует ли степень тройки, оканчивающаяся на 0001 ?

\problem{8}
Докажите, что $2^{3^k}+1$ делится на $3^{k+1}$.

\problem{9} Пусть $m$ и $n$ натуральные числа.

a) Докажите, что $\varphi(n) \cdot \varphi(m)=\varphi($ НОД $(m, n)) \cdot \varphi($ НОК $(m, n))$.

б) Докажите \textbf{усиление теоремы Эйлера}. Положим $L(m)=\operatorname{HOK}\left(\varphi\left(p_1^{t_1}\right), \ldots, \varphi\left(p_k^{t_k}\right)\right)$. Тогда, если $(a, m)=1$, то $a^{L(m)}-1$ делится на $m$.
